% ═══════════════════════════════════════════════════════════════
% LEHRERHINWEISE LH-03a: Las habitaciones
% Spanisch Klasse 7 - Sequenz 3: Mi casa
% ═══════════════════════════════════════════════════════════════
% Fachfarbe: #c41e3a (Karminrot)
%
% Enthält:
% - Kurzplan (Stundenverlauf)
% - Differenzierungshinweise [B]/[S]/[E]
% - Häufige Fehler + Reaktionen
% - Material-Checkliste
% - Lösungen
% ═══════════════════════════════════════════════════════════════

\documentclass[11pt, a4paper]{article}

% ═══════════════════════════════════════════════════════════════
% SW-MODUS TOGGLE
% ═══════════════════════════════════════════════════════════════
\newif\ifswmodus
\swmodusfalse

% ═══════════════════════════════════════════════════════════════
% PAKETE
% ═══════════════════════════════════════════════════════════════

\usepackage[utf8]{inputenc}
\usepackage[T1]{fontenc}
\usepackage[ngerman]{babel}
\usepackage{helvet}
\renewcommand{\familydefault}{\sfdefault}

\usepackage[margin=2cm]{geometry}
\usepackage{enumitem}
\usepackage{tabularx}
\usepackage{booktabs}
\usepackage{longtable}

\usepackage{xcolor}
\usepackage{amssymb}
\usepackage{tcolorbox}
\tcbuselibrary{skins}

\usepackage{fancyhdr}
\usepackage{lastpage}
\usepackage{needspace}

% ═══════════════════════════════════════════════════════════════
% FARBEN
% ═══════════════════════════════════════════════════════════════

\definecolor{spanisch-color}{HTML}{c41e3a}
\definecolor{grau-color}{HTML}{888888}
\definecolor{hellgrau-color}{HTML}{f5f5f5}
\definecolor{basis-color}{HTML}{2a7a4b}
\definecolor{standard-color}{HTML}{2c5aa0}
\definecolor{erweiterung-color}{HTML}{7b1fa2}
\definecolor{loesung-color}{HTML}{c62828}

\definecolor{sw-dark}{HTML}{000000}
\definecolor{sw-gray}{HTML}{666666}
\definecolor{sw-white}{HTML}{ffffff}

\ifswmodus
    \colorlet{spanisch}{sw-dark}
    \colorlet{grau}{sw-gray}
    \colorlet{hellgrau}{sw-white}
    \colorlet{basis}{sw-dark}
    \colorlet{standard}{sw-dark}
    \colorlet{erweiterung}{sw-dark}
    \colorlet{loesung}{sw-dark}
\else
    \colorlet{spanisch}{spanisch-color}
    \colorlet{grau}{grau-color}
    \colorlet{hellgrau}{hellgrau-color}
    \colorlet{basis}{basis-color}
    \colorlet{standard}{standard-color}
    \colorlet{erweiterung}{erweiterung-color}
    \colorlet{loesung}{loesung-color}
\fi

\newcommand{\fachfarbe}{spanisch}

% ═══════════════════════════════════════════════════════════════
% VARIABLEN
% ═══════════════════════════════════════════════════════════════

\newcommand{\thema}{Las habitaciones}
\newcommand{\sequenz}{03}
\newcommand{\block}{a}
\newcommand{\fach}{Spanisch}
\newcommand{\klasse}{7}
\newcommand{\dauer}{90 Min. (Doppelstunde)}
\newcommand{\lerngruppengroesse}{24}
\newcommand{\datum}{\today}

% ═══════════════════════════════════════════════════════════════
% KOPF- UND FUSSZEILE
% ═══════════════════════════════════════════════════════════════

\pagestyle{fancy}
\fancyhf{}
\renewcommand{\headrulewidth}{0.4pt}
\renewcommand{\footrulewidth}{0.4pt}

\lhead{\fach{} -- Klasse \klasse}
\chead{\textbf{LH-\sequenz\block{} (Lehrerhinweise)}}
\rhead{\datum}

\lfoot{}
\cfoot{}
\rfoot{Seite \thepage\ von \pageref{LastPage}}

% ═══════════════════════════════════════════════════════════════
% BEFEHLE
% ═══════════════════════════════════════════════════════════════

\newcommand{\B}{\textcolor{basis}{\textbf{[B]}}}
\newcommand{\Std}{\textcolor{standard}{\textbf{[S]}}}
\newcommand{\E}{\textcolor{erweiterung}{\textbf{[E]}}}

\newcommand{\loesung}[1]{\textcolor{loesung}{\textbf{#1}}}

\newcommand{\es}[1]{\textcolor{\fachfarbe}{\textbf{#1}}}

\newtcolorbox{kurzplan}{
    colback=hellgrau,
    colframe=\fachfarbe,
    colbacktitle=white,
    coltitle=\fachfarbe,
    boxrule=1pt,
    arc=0pt,
    fonttitle=\bfseries,
    attach boxed title to top left={yshift=-2mm, xshift=5mm},
    boxed title style={boxrule=0pt, colback=white},
    title=Kurzplan
}

\newtcolorbox{differenzierung}{
    colback=white,
    colframe=grau,
    colbacktitle=white,
    coltitle=grau,
    boxrule=0.5pt,
    arc=0pt,
    fonttitle=\bfseries,
    attach boxed title to top left={yshift=-2mm, xshift=5mm},
    boxed title style={boxrule=0pt, colback=white},
    title=Differenzierung
}

\newtcolorbox{fehlerbox}{
    colback=white,
    colframe=loesung,
    colbacktitle=white,
    coltitle=loesung,
    boxrule=0.5pt,
    arc=0pt,
    fonttitle=\bfseries,
    attach boxed title to top left={yshift=-2mm, xshift=5mm},
    boxed title style={boxrule=0pt, colback=white},
    title=Häufige Fehler
}

\newtcolorbox{materialbox}{
    colback=white,
    colframe=\fachfarbe,
    colbacktitle=white,
    coltitle=\fachfarbe,
    boxrule=0.5pt,
    arc=0pt,
    fonttitle=\bfseries,
    attach boxed title to top left={yshift=-2mm, xshift=5mm},
    boxed title style={boxrule=0pt, colback=white},
    title=Material-Checkliste
}

% ═══════════════════════════════════════════════════════════════
% DOKUMENT
% ═══════════════════════════════════════════════════════════════

\begin{document}

% ═══════════════════════════════════════════════════════════════
% KOPFBEREICH
% ═══════════════════════════════════════════════════════════════

\begin{center}
    {\Large\bfseries\textcolor{\fachfarbe}{Lehrerhinweise: \thema}}\\[0.3em]
    {\normalsize LH-\sequenz\block{} | \dauer}
\end{center}

\vspace{10pt}

% ═══════════════════════════════════════════════════════════════
% 1. KURZPLAN
% ═══════════════════════════════════════════════════════════════

\section*{1. Stundenverlauf (Kurzplan)}

\textbf{1. Stunde (45 Min.)}

\begin{tabularx}{\textwidth}{|c|l|X|l|}
\hline
\textbf{Zeit} & \textbf{Phase} & \textbf{Inhalt} & \textbf{Material} \\
\hline
0--5 & Einstieg & Bild span. Haus zeigen, ``¿Qué ves?'' & Projektor \\
\hline
5--15 & Input & 8 Räume mit Bildkarten + Audio einführen & Karten, Audio \\
\hline
15--25 & Übung 1 & Chorisches Sprechen, Genus-Regel erklären & -- \\
\hline
25--35 & Übung 2 & Bild-Wort-Memory in Gruppen & Memory-Karten \\
\hline
35--45 & Sicherung & Merkkasten auf AB ausfüllen & AB-03a \\
\hline
\end{tabularx}

\vspace{1em}

\textbf{2. Stunde (45 Min.)}

\begin{tabularx}{\textwidth}{|c|l|X|l|}
\hline
\textbf{Zeit} & \textbf{Phase} & \textbf{Inhalt} & \textbf{Material} \\
\hline
0--5 & Aktivierung & Quiz: ``¿Cómo se dice Küche?'' & -- \\
\hline
5--15 & Vertiefung & AB Aufgabe 1-2 (Zuordnung, Artikel) & AB-03a \\
\hline
15--25 & Anwendung & Hausplan beschriften, Partnercheck & AB-03a \\
\hline
25--35 & Dialog & ``¿Dónde está...?'' -- ``Está en...'' üben & -- \\
\hline
35--40 & Erweiterung & \E{}: Kulturvergleich Wohnen Spanien & Tafel \\
\hline
40--45 & Abschluss & Zusammenfassung, HA erklären & -- \\
\hline
\end{tabularx}

% ═══════════════════════════════════════════════════════════════
% 2. DIFFERENZIERUNG
% ═══════════════════════════════════════════════════════════════

\section*{2. Differenzierung}

\begin{differenzierung}
\textbf{Legende:} \B{} Basis (BR) | \Std{} Standard (MR) | \E{} Erweiterung (GY)

\vspace{0.5em}

\textbf{WICHTIG:} Diese Marker sind NUR für die Lehrkraft. Auf Schülermaterial erscheinen KEINE Niveaumarkierungen!
\end{differenzierung}

\vspace{1em}

\begin{tabularx}{\textwidth}{|l|X|X|X|}
\hline
& \B{} \textbf{Basis} & \Std{} \textbf{Standard} & \E{} \textbf{Erweiterung} \\
\hline
\textbf{Aufgabe 1} & Bild-Wort-Zuordnung mit Wortliste & Zuordnung selbstständig & + Sätze mit Raum bilden \\
\hline
\textbf{Aufgabe 2} & Artikel mit Genus-Regel Hilfe & Artikel selbstständig & + Begründung nennen \\
\hline
\textbf{Aufgabe 3} & Hausplan mit 4 Räumen & Hausplan mit 6 Räumen & + Eigenes Haus zeichnen \\
\hline
\textbf{Aufgabe 4} & Dialog nachsprechen & Dialog mit Stichworten & Freier Dialog \\
\hline
\end{tabularx}

\vspace{1em}

\textbf{Konkrete Hinweise:}
\begin{itemize}
    \item \B{} SuS mit LRS: Zeitzugabe (+10 Min.), Wortliste mit Bildern
    \item \B{} DaZ-SuS: Raumbezeichnungen auch auf Deutsch vorab klären
    \item \E{} Leistungsstarke: Eigenen Hausplan erstellen und präsentieren
    \item \E{} Zusatz: Kulturvergleich Wohnen in Spanien vs. Deutschland recherchieren
\end{itemize}

% ═══════════════════════════════════════════════════════════════
% 3. HÄUFIGE FEHLER
% ═══════════════════════════════════════════════════════════════

\section*{3. Häufige Fehler}

\begin{fehlerbox}
\begin{tabularx}{\textwidth}{|l|X|X|}
\hline
\textbf{Fehler} & \textbf{Beispiel} & \textbf{Reaktion} \\
\hline
Falsches Genus & *\es{el cocina} statt \es{la cocina} & Farbcodierung (el=blau, la=rot), Endungsregel \\
\hline
Aussprache \textit{j} & falsche Aussprache & ``ch'' wie in ``Buch'', chorisch üben \\
\hline
Verwechslung salón/comedor & Unsicherheit Wohn-/Esszimmer & Visualisierung: Sofa vs. Esstisch \\
\hline
*el terraza & Endung -a ignoriert & Regel: -a $\rightarrow$ la (mit Ausnahmen) \\
\hline
¿Dónde...? unvollständig & *``¿Dónde tu madre?'' & Immer \es{¿Dónde está...?} einfordern \\
\hline
\end{tabularx}
\end{fehlerbox}

% ═══════════════════════════════════════════════════════════════
% 4. MATERIAL-CHECKLISTE
% ═══════════════════════════════════════════════════════════════

\section*{4. Material-Checkliste}

\begin{materialbox}
\begin{itemize}[label=$\square$]
    \item AB-03a.pdf (\lerngruppengroesse{} Kopien)
    \item ML-03a.pdf (1$\times$ für Lehrkraft)
    \item Bildkarten: 8 Räume (laminiert, Klassensatz)
    \item Memory-Karten: 8 Paare (6 Sets für Gruppen)
    \item Projektor + Bild: Spanisches Haus (z.B. Andalusien)
    \item Lehrwerk: Qué pasa 1, S. 36-39
    \item Audio-Track (falls Lehrwerk-CD verwendet)
    \item Tafelmarker (blau für \es{el}, rot für \es{la})
    \item Optional: LP-03a.html auf Tablets
\end{itemize}
\end{materialbox}

% ═══════════════════════════════════════════════════════════════
% 5. LÖSUNGEN
% ═══════════════════════════════════════════════════════════════

\section*{5. Lösungen AB-03a}

\textbf{Merkkasten:}
\begin{tabular}{llll}
\es{la cocina} & = die Küche & \es{el salón} & = das Wohnzimmer \\
\es{el dormitorio} & = das Schlafzimmer & \es{el baño} & = das Badezimmer \\
\es{el comedor} & = das Esszimmer & \es{el jardín} & = der Garten \\
\es{el garaje} & = die Garage & \es{la terraza} & = die Terrasse \\
\end{tabular}

\vspace{1em}

\textbf{Aufgabe 1: Zuordnung Bild-Wort} (4 Punkte)

\begin{tabular}{ll}
Bild A (Herd) & $\rightarrow$ \loesung{la cocina} \\
Bild B (Sofa) & $\rightarrow$ \loesung{el salón} \\
Bild C (Bett) & $\rightarrow$ \loesung{el dormitorio} \\
Bild D (Badewanne) & $\rightarrow$ \loesung{el baño} \\
\end{tabular}

\vspace{1em}

\textbf{Aufgabe 2: Artikel einsetzen} (4 Punkte)

\begin{tabular}{ll}
a) & \loesung{la} cocina \\
b) & \loesung{el} jardín \\
c) & \loesung{la} terraza \\
d) & \loesung{el} garaje \\
\end{tabular}

\vspace{1em}

\textbf{Aufgabe 3: Hausplan beschriften} (6 Punkte)

\begin{enumerate}
    \item \loesung{el salón} (oben links)
    \item \loesung{la cocina} (oben rechts)
    \item \loesung{el dormitorio} (unten links)
    \item \loesung{el baño} (unten rechts)
    \item \loesung{el jardín} (außen rechts)
    \item \loesung{el comedor} (Dachbereich)
\end{enumerate}

\textit{Bewertung: Je 1P pro korrekter Beschriftung inkl. Artikel}

\vspace{1em}

\textbf{Aufgabe 4: Minidialog} (6 Punkte)

Beispiellösung:
\begin{itemize}
    \item A: ¿Dónde está tu madre? -- B: Está en la cocina.
    \item A: ¿Y tu padre? -- B: Está en el jardín.
    \item A: ¿Dónde duermes tú? -- B: Duermo en el dormitorio.
\end{itemize}

\textit{Bewertung: Je 2P pro sinnvollen Austausch (Frage + Antwort). Variationen akzeptiert.}

% ═══════════════════════════════════════════════════════════════
% 6. HAUSAUFGABE
% ═══════════════════════════════════════════════════════════════

\section*{6. Hausaufgabe}

\begin{itemize}
    \item Vokabeln lernen: 8 Räume mit Artikeln
    \item AB-03a fertigstellen (falls nicht abgeschlossen)
    \item Optional: Lernpfad LP-03a.html bearbeiten
    \item Optional (\E{}): Eigenes Haus/Wohnung zeichnen und beschriften
\end{itemize}

% ═══════════════════════════════════════════════════════════════
% 7. KULTURINFO (für Lehrkraft)
% ═══════════════════════════════════════════════════════════════

\section*{7. Kulturinfo: Wohnen in Spanien}

\textbf{Für Kulturvergleich (\E{}) oder Ausblick:}

\begin{itemize}
    \item \textbf{Piso vs. Casa:} In Städten dominieren Wohnungen (piso), auf dem Land Häuser (casa)
    \item \textbf{Patio:} Innenhof als zentraler Lebensraum, besonders in Andalusien
    \item \textbf{Terraza:} Dachterrasse oder Balkon -- wichtig für Soziales
    \item \textbf{Persianas:} Rollläden wegen Hitze und Siesta
    \item \textbf{Azulejos:} Traditionelle Kacheln, oft in Küche und Bad
\end{itemize}

% ═══════════════════════════════════════════════════════════════
% 8. REFLEXIONSNOTIZEN
% ═══════════════════════════════════════════════════════════════

\section*{8. Reflexionsnotizen (nach Durchführung)}

\begin{tabularx}{\textwidth}{|l|X|}
\hline
\textbf{Was lief gut?} & \\[2em]
\hline
\textbf{Was lief nicht?} & \\[2em]
\hline
\textbf{Zeitmanagement} & \\[2em]
\hline
\textbf{Für nächstes Mal} & \\[2em]
\hline
\end{tabularx}

\end{document}
